% ---
% title: Rethink Diffusion Models for Feature Modeling
% date: 2026-06-21
% tag: Notes
% excerpt: 把扩散模型当作特征分布建模器，讨论小样本和高维特征两个场景：什么时候有用，什么时候只是在学伪分布。
% ---
扩散模型在图像、音频和视频生成上已经证明有效：前向过程逐步加噪，反向过程逐步去噪，模型学到的是数据分布。

把同样的过程搬到特征空间，问题就不一样了。这里要建模的不是图像，而是特征向量。特征可以是深度网络中的 embedding，也可以是某个任务里的表征向量：

\[
z \in \mathbb{R}^{d}
\]

特征空间里要建模的是 \(p(z)\)，不是原始数据分布 \(p(x)\)。变量从 \(x\) 换成 \(z\)，问题却变了：特征空间更抽象、维度更高、更稀疏，也更依赖任务语义。

下面看两个场景：小样本特征建模和高维特征建模。两个场景都绕不开同一个问题：扩散模型的分布建模能力，能不能抵消小样本和高维空间带来的统计困难。

\section{扩散模型作为特征分布建模器}

给定特征向量 \(z\)，扩散模型学的是它的分布：

\[
p(z)
\]

扩散模型先通过前向过程向特征中逐步加入噪声：

\[
q(z_t \mid z_{t-1})
\]

然后学习一个反向过程，从噪声特征中逐步恢复原始特征：

\[
p_\theta(z_{t-1} \mid z_t)
\]

常见的加噪形式可以写作：

\[
z_t = \sqrt{\alpha_t}z_0 + \sqrt{1-\alpha_t}\epsilon
\]

其中，\(z_0\) 是原始特征，\(\epsilon\) 是高斯噪声。

训练目标通常是让模型预测噪声：

\[
\mathcal{L}_{diff} =
\mathbb{E}_{z_0,t,\epsilon}
\left[
\|\epsilon - \epsilon_{\theta}(z_t,t)\|^2
\right]
\]

这样学出来的是特征空间中的局部结构、密度形态和去噪方向，分类边界不在目标里。

吸引力也在这里：它要学的是特征分布本身，判别结果不是终点。

\section{场景一：小样本特征建模}

小样本指每类只有少量样本，却仍要尽量准确地估计该类特征的分布。例如：

\[
{z_1,z_2,\dots,z_n}, \quad n \ll d
\]

有时 \(n\) 未必小于 \(d\)，但相对于特征空间的复杂度，样本仍然不够。

困难很直接：数据不足以支撑复杂分布的估计。

\subsection{扩散模型在小样本场景中的潜在优势}

扩散模型可以通过去噪任务构造大量扰动样本。

对于每个原始特征 \(z_0\)，可以在不同时间步生成不同噪声版本：

\[
z_t = \sqrt{\alpha_t}z_0 + \sqrt{1-\alpha_t}\epsilon
\]

这样，一个样本可以对应多个噪声状态：

\[
z_0 \rightarrow {z_t}\big |_{t=1}^{T}
\]

这在表面上增加了训练信号。

换个角度看，扩散训练要求模型从邻域里的噪声点回到数据流形附近，只记住原始样本点不够。

在小样本下，噪声扰动提供了局部训练信号，模型学到的是样本附近的几何结构，也能捕捉维度之间的高阶相关性。生成的特征可以用于数据增强，并在特征空间里起到平滑正则的作用。有限的样本因此被扩展成一组去噪任务。

如果真实分布足够平滑、样本之间共享结构，扩散模型可能比直接记忆样本更有效。

\subsection{小样本场景中的主要风险：伪分布学习}

小样本也是扩散模型最容易出问题的场景之一。模型参数量通常较大，小样本提供的信息有限。模型可能只学到训练样本附近的局部邻域，总体分布 \(p(z)\) 并没有被学到：

\[
p_\theta(z) \approx p_{\text{train-neighborhood}}(z)
\]

模型可能在训练样本周围生成看似合理的特征，这些特征并不代表真实总体分布。

这就是伪分布建模：模型看起来学到了分布，实际只是把少数样本平滑地记了下来。

生成结果可能比原始样本更多、更平滑，却不一定更真实。

\subsection{小样本下的数据增强不一定可靠}

使用扩散模型生成特征样本是一种自然想法：

\[
\tilde{z} \sim p_\theta(z)
\]

然后将生成特征加入下游训练。

问题在于，\(p_\theta(z)\) 本身就是从少量样本估出来的，生成样本很可能只是训练样本附近的插值。这类样本能让训练更稳定，却未必提高泛化。样本差异很大时，小样本扩散模型还可能把不该连的区域连起来。

例如，两个样本 \(z_a\) 和 \(z_b\) 可能属于同一类，但它们在特征空间中处于两个不同模式：

\[
z_a \in \mathcal{M}_1,\quad z_b \in \mathcal{M}_2
\]

如果扩散模型强行在二者之间形成连续路径，就可能生成不真实的中间特征。

小样本扩散建模要问的是：生成的特征是否落在真实语义流形上。数量本身说明不了问题。

\subsection{小样本下更合理的使用方式}

小样本下，扩散模型不适合当完整的分布估计器。更合适的是让它做局部建模或正则化。

例如，可以只让扩散模型学习局部扰动恢复：

\[
z_t \rightarrow z_0
\]

不要求它完整建模整个类别分布。

也可以使用条件信息限制生成范围：

\[
p_\theta(z \mid c)
\]

其中，\(c\) 可以是类别、原型、文本描述、聚类中心或其他先验信息。

还可以将扩散模型用于原型附近的局部扩增：

\[
z = \mu_c + \delta
\]

其中，\(\mu_c\) 是类别原型，\(\delta\) 是局部变化。

这样模型学的是围绕少量原型的局部变化，而不是任意高维分布。

对于小样本场景，更稳妥的思路是：

\[
\text{Few samples}
\rightarrow
\text{Prototype or prior}
\rightarrow
\text{Local diffusion}
\rightarrow
\text{Regularized feature augmentation}
\]

扩散模型应该依赖先验，单靠少量数据学复杂分布并不可靠。

\section{场景二：高维特征建模}

高维特征建模更常见。

现代深度模型产生的 embedding 往往维度很高：

\[
z \in \mathbb{R}^{512},\quad
z \in \mathbb{R}^{768},\quad
z \in \mathbb{R}^{1024},\quad
z \in \mathbb{R}^{4096}
\]

在这些空间中，扩散模型需要学习的是高维联合分布：

\[
p(z_1,z_2,\dots,z_d)
\]

高维不只是维度更多。维度一高，统计估计、几何结构和优化稳定性都会出问题。

\subsection{高维空间中的数据稀疏性}

高维空间里样本非常稀疏。样本数量看着不少，相对于整个空间的体积仍然极少。

维度 \(d\) 增大时，覆盖特征空间所需的样本量增长得很快。

扩散模型很难可靠估计完整的高维分布：

\[
p(z)
\]

它更可能只学习到训练样本附近的一些局部结构。

生成样本可能缺乏多样性，去噪方向在分布外区域也可能不稳定。

在高维空间中，一个噪声扰动可能很容易把特征推出真实流形：

\[
z_t \notin \mathcal{M}
\]

其中，\(\mathcal{M}\) 表示真实特征流形。如果大部分训练都发生在流形外，模型学到的去噪路径有没有语义意义，就需要重新检查。

\subsection{高维特征通常位于低维流形上}

特征向量维度很高，但往往并不均匀分布在整个空间里。

更合理的假设是，真实特征集中在某个低维流形附近：

\[
z \in \mathcal{M}, \quad \dim(\mathcal{M}) \ll d
\]

要建模的是嵌在 \(\mathbb{R}^{d}\) 里的低维结构，而不是整个空间。

如果扩散模型直接在原始高维空间中加噪，样本会被从流形上推到流形外。

矛盾在于：扩散模型靠加噪学习恢复路径，但高维噪声会制造大量不属于真实数据结构的点。

模型的能力可能花在从流形外的噪声回到流形上，流形内部真正有价值的变化反而没学到。

\subsection{欧氏噪声不一定符合特征语义}

标准扩散过程通常假设高斯噪声是合理的：

\[
\epsilon \sim \mathcal{N}(0,I)
\]

并使用：

\[
z_t = \sqrt{\alpha_t}z_0 + \sqrt{1-\alpha_t}\epsilon
\]

但在高维特征空间中，欧氏距离和语义距离未必一致。

两个特征在欧氏距离上很近，语义未必相近；距离很远，语义也未必差得远。

因此，直接在特征空间加高斯噪声可能会破坏原有语义结构。

更一般地说，特征空间可能具有非欧氏几何：

\[
d_{\text{semantic}}(z_i,z_j) \neq \|z_i-z_j\|_2
\]

如果加噪依赖欧氏距离，模型学到的更接近数值空间里的去噪，语义空间里的生成并没有被直接建模。

\subsection{高维噪声预测的优化偏差}

扩散模型通常优化噪声预测误差：

\[
\mathcal{L}_{diff}=
\mathbb{E}
\left[
\|\epsilon - \epsilon_\theta(z_t,t)\|^2
\right]
\]

高维特征里，不同维度的重要性并不相同。某些维度数值方差大，语义意义却弱；某些维度方差小，对下游任务却很关键。

直接用均方误差，模型会偏向整体的数值重建，任务相关的结构容易被忽略。

最终学到的可能是一个平均意义上合理的去噪器，关键判别信息或细粒度结构未必保得住。

高维特征扩散要回答的问题就变成了：优化目标真的对应我们想要的特征质量吗？

\section{小样本与高维特征的共同矛盾}

小样本和高维经常同时出现。比如只有少量样本，每个样本的特征维度却很高：

\[
n \ll d
\]

这是最困难的情况之一。

扩散模型同时受两头挤压：小样本让分布估计不可靠，高维让数据更加稀疏。

表达能力强不等于能解决统计不充分的问题。没有足够先验和约束，强模型反而更容易过拟合少量样本。

小样本高维场景里，直接训练更大的扩散模型不是办法。应该先把建模难度降下来。

\section{可能的改进方向}

\subsection{先降维，再扩散}

直接在高维特征上扩散通常不是最稳妥的选择。

可以先通过编码器或投影层得到低维表示：

\[
h = g_\phi(z)
\]

其中：

\[
h \in \mathbb{R}^{k}, \quad k \ll d
\]

然后在低维空间中建模：

\[
p_\theta(h)
\]

这样可以减少数据稀疏性和优化不稳定性。

降维不能只看压缩率，结构信息也要保留。压得太狠，扩散模型学到的空间更简单，关键变化也一起丢了。

\subsection{使用原型或先验约束}

在小样本场景中，可以引入原型：

\[
\mu_c = \frac{1}{n_c}\sum_{i=1}^{n_c} z_i
\]

然后建模样本相对于原型的局部偏移：

\[
\delta_i = z_i - \mu_c
\]

扩散模型可以只建模：

\[
p_\theta(\delta \mid c)
\]

而不是完整建模：

\[
p_\theta(z \mid c)
\]

这样任务难度降低了，模型关注的是类内变化，不必从少量样本里估计完整分布。

\subsection{使用条件扩散减少不确定性}

条件扩散可以将建模目标从无条件分布：

\[
p(z)
\]

转化为条件分布：

\[
p(z \mid c)
\]

条件 \(c\) 可以是类别、属性、原型、上下文、聚类编号或其他先验。

条件越明确，模型需要解释的不确定性越小。

高维和小样本场景尤其如此。无条件建模要模型同时解释多种变化来源，条件建模把问题拆小了。

\subsection{使用流形感知的噪声过程}

如果特征分布位于低维流形上，那么加噪过程也应该尽量尊重这个流形结构。

标准高斯噪声：

\[
\epsilon \sim \mathcal{N}(0,I)
\]

可能不是最合适的选择。

更合理的方向是使用流形感知扰动，例如：

\[
z_t = z_0 + \epsilon_{\mathcal{T}}
\]

其中，\(\epsilon_{\mathcal{T}}\) 表示沿流形切空间的扰动。

这类方法在更符合数据结构的方向上制造扰动，不把样本推向任意高维方向。

实现会更复杂，但原则很清楚：特征扩散不该照搬像素空间的加噪方式。

\subsection{将扩散目标与任务目标结合}

如果特征建模最终服务于某个任务，那么单纯优化去噪误差可能不够。

可以将扩散损失与任务损失结合：

\[
\mathcal{L}
=
\mathcal{L}_{\text{diff}}
+
\lambda \mathcal{L}_{\text{task}}
\]

其中，\(\mathcal{L}_{task}\) 可以是分类、检索、度量学习、对比学习或其他任务相关目标。

这样扩散模型在数值重建之外，也会保住任务相关结构。

\section{扩散模型在特征建模中的合理定位}

在小样本和高维场景中，扩散模型不应该被看作万能的分布学习器。

它更像一个受约束的局部生成器，或者特征空间里的正则化器。

适合做的是：在样本邻域内生成合理扰动，学习特征流形附近的去噪方向，对特征做补全、修复或平滑，在低维 latent 中建模结构变化，以及在条件约束下生成特征增强样本。

不适合直接做的是：从极少样本里学完整的高维分布，在原始高维 embedding 上无约束生成，用欧氏高斯噪声替代语义扰动，用去噪误差直接代表特征质量，以及把生成样本默认当成有效的数据增强。

更合理的设计不是：

\[
\text{Raw high-dimensional features}
\rightarrow
\text{Unconditional diffusion}
\]

而是：

\[
\text{Raw features}
\rightarrow
\text{Compact task-aware representation}
\rightarrow
\text{Conditional or local diffusion}
\rightarrow
\text{Feature modeling / augmentation / regularization}
\]

扩散模型用于特征建模，决定效果的是建模空间是否合理，模型容量本身说明不了太多。

更值得建模的往往是经过压缩、对齐和约束之后的特征子空间，而不是原始高维 embedding。

小样本和高维场景里，扩散模型要依赖低维表示、条件约束和任务先验。缺了这些，生成能力很容易退化成对稀疏训练样本的平滑记忆。
